\section{Independent exact-degree certificates for upper centering}
\label{sec:upper-centered-degrees}

This section concerns the fixed-blank, ten-positive-variable circuit with
$X=wq$, $q=\beta J+1$, and the unchanged compact44 coefficient tuple.
The coefficient ports are constants when polynomial degree is measured.
Write $A=\mathrm{Am1}+1$, $H=4A-5$, $P_1=q^3(q-1)$ and
$M=P_1\mu^2H$.  The registers retained from the source satisfy
\begin{align*}
 D_p&=A^2-1,&D_f&=D_p I^2+1,&Q&=D_p^2I^2,&\delta&=Q-1,\\
 V&=D_f(c(T+\alpha)-R)-c,&D_{\rm num}&=QV^2-1,\\
 C_f&=(A+1)I,&C_t&=(V_0+1)k,&
 D_\tau&=V_0(V_0+1)k^2+1.
\end{align*}
Let $L=k(T+\alpha)$ and
\[
 E=n-LB_t+\frac{K_{\rm pos}}H
   +\frac{(q-1)Y-q^3B_t}{P_1}-\frac e{\mu^2}.
\]
The strong-center circuit has $Z_{\rm num}=MZ$ with
$Z=V+C_f+E$.  The two-center circuit instead uses
$Z=V+C_f+3C_t+E$.  These are rational identities between literal
register expressions; $Z_{\rm num}$ itself is a polynomial because the
displayed denominators are cleared by $M$.

\paragraph{Theorem D.1.}
The strong-center polynomial has exact total degree $358$.  The
two-center polynomial has exact total degree $354$.  Their top
homogeneous forms are, respectively,
\begin{align*}
 [N_f]_{358}
   &=-9\,2^{22}(\beta J)^{188}w^{36}Y^{74}i^{18}
       (\eta+\zeta)^{38}(\beta J-F)^4,\\
 [N_{ft}]_{354}
   &=3\,2^{23}(\beta J)^{183}w^{35}Y^{74}i^{19}
       (\eta+\zeta)^{39}(\beta J-F)^4.
\end{align*}
In particular these forms are nonzero for every fixed $\beta>0$.
Neither equality is inferred from a syntactic maximum-degree propagation,
which overlooks the cancellation in the final norm.

\paragraph{Formal setup.}
Replace each of the ten witnesses by $t$ times a separate indeterminate,
and work in the Laurent-series field $\mathcal K((t^{-1}))$, where
$\mathcal K$ is the rational function field in those ten indeterminates.
The fixed integer coefficient ports are in the coefficient field.
For a nonzero Laurent series, ``degree'' means its largest exponent of
$t$.  For polynomials this agrees with the ordinary total degree before
the scaling.  A rational function with square leading coefficient has a
formal square root after choosing its sign; here all three needed leading
coefficients are explicit squares.  The usual binomial series proves
existence, because each relative correction has strictly negative degree.
No analytic convergence or assertion about positive solution values is
used.

Choose formal roots
\[
 f^2=D_f,\qquad \tau^2=D_\tau,\qquad
 g^2=D_{\rm num}/\delta
\]
with $f$ near $AI$, $\tau$ near $C_t$, and $g$ near $V$.
The following exact identities control all cancellations:
\begin{align*}
 (AI-f)(AI+f)&=I^2-1,\\
 (C_t-\tau)(C_t+\tau)&=kC_t-1,\\
 (g-V)(g+V)&=(V^2-1)/\delta.
\end{align*}
Thus $\deg(AI-f)=0$, $\deg(C_t-\tau)=1$, and
$\deg(g-V)=-6$.  In particular
$C_f-f=I+(AI-f)$ has degree $5$, with top form equal to that of $I$.

For completeness, the register degrees relevant to this argument are
\[
\begin{array}{c|rrrrrrrrrrrr}
\text{register}&A&I&R&V&C_f&C_t&\delta&M&E&f&\tau&g\\\hline
\text{degree}&5&5&4&24&10&9&30&13&\leq4&10&9&24.
\end{array}
\]
The bound on $E$ follows directly from $\deg(LB_t)=4$,
$\deg(K_{\rm pos}/H)\leq-1$,
$\deg(((q-1)Y-q^3B_t)/P_1)\leq1$, and
$\deg(e/\mu^2)\leq0$.  The term $c(T+\alpha)$ has degree $3$ whereas
$R$ has degree $4$, so the leading form of $V$ is $-[D_f]_{20}[R]_4$.

\paragraph{Norm factorization.}
The literal final circuit computes
\[
 \bigl((s+c_0-a-b)^2-4(sc_0+ab)\bigr)^2-64sc_0ab,
\]
where
$s=\delta Z_{\rm num}^2$, $c_0=M^2D_{\rm num}$,
$a=9\delta M^2D_\tau$, and $b=\delta M^2D_f$.
Successive multiplication of conjugate pairs gives the exact identity
\[
 N=\delta^4M^8
       \prod_{\sigma,\rho,\epsilon\in\{-1,1\}}
           (Z+\sigma g+3\rho\tau+\epsilon f).
\]
This is an algebraic identity; the formal choices of signs of the three
square roots merely permute its eight factors.

Put $q_0=\beta J$, $k_0=\eta+\zeta$, and abbreviate the following
nonzero homogeneous forms:
\begin{align*}
 a_0&=q_0^3wY,&j_0&=ik_0^2Y^2,&
 r_0&=q_0^3(q_0-F),\\
 t_0&=q_0^5wY^2k_0,&f_0&=a_0j_0,&
 v_0&=-a_0^2j_0^2r_0,\\
 d_0&=a_0^4j_0^2,&m_0&=4q_0^8a_0.
\end{align*}
They are the top forms of $A,I,R,\tau,f,V,\delta,M$, in the
corresponding order.  The four factors with $\sigma=1$ all have degree
$24$ and top form $2v_0$.  Among the factors with $\sigma=-1$, those
with $\epsilon=1$ both have degree $10$ and top form $2f_0$.

For strong centering, the last two factors have degree $9$ and top forms
$3t_0$ and $-3t_0$.  The total degree is therefore
\[
 4(30)+8(13)+4(24)+2(10)+2(9)=358,
\]
and the product of top forms is
$-576d_0^4m_0^8v_0^4f_0^2t_0^2$.
Substitution of the displayed abbreviations gives the first formula in
the theorem.

For two-center centering, the last two factors instead have degrees
$9$ and $5$, with top forms $6t_0$ and $j_0$.  The latter conclusion
uses $\deg(C_f-f)=5$, while the other three terms $E$,
$3(C_t-\tau)$, and $V-g$ have degrees at most $4,1,-6$.
Thus
\[
 4(30)+8(13)+4(24)+2(10)+9+5=354,
\]
and the top-form product is
$384d_0^4m_0^8v_0^4f_0^2t_0j_0$.
This gives the second formula.  Each displayed top form is a nonzero
polynomial, so the argument proves both the upper degree bound and its
attainment. \hfill$\square$

\paragraph{Independent exact ray check.}
The accompanying standard-library Python program independently rebuilds
both transformations from the pinned source.  It evaluates every circuit
gate as a univariate integer polynomial after putting every witness
equal to $t$ and setting $(\beta,K_0,M_C,M_F)=(2,3,5,7)$.
The resulting leading terms are exactly
$-9\,2^{248}t^{358}$ and $3\,2^{245}t^{354}$.
These finite computations check the formulas on one exact ray; the
Laurent-series argument above supplies the global bounds and proves
nonvanishing for the actual fixed coefficient tuple.

The identity $C_f=(A+1)I$, $S=(A-1)C_f=(A^2-1)I$ also permits the
shared evaluation $D_f=SI+1$, $Q=S^2$.  Deleting the old separate
$D_p$ and $I^2$ gates saves one operation without changing any
polynomial, and therefore does not change either degree certificate.

\subsection{The direct-$X$ siblings}

For the direct source, $X=q+x$, the carrier is
$B=(qK_0+X)C-qF$, and
$P_1=q^4(q-1)$.  In particular its extra factor of $q$ must be retained
when computing the degree of the norm prefactor.  Here $M=P_1\mu^2H$
still has degree $13$, while $\delta$ has degree $26$.
The centered rational identity uses
\[
 E=n-LB+\frac{K_{\rm pos}}H+\frac{Y}{q^3}
       -\frac{B}{q(q-1)}-\frac{X}{q^4}-\frac e{\mu^2}.
\]
Its degree is at most $4$.  The relevant degrees become
\[
\begin{array}{c|rrrrrrrrrrrr}
\text{register}&A&I&R&V&C_f&C_t&\delta&M&E&f&\tau&g\\\hline
\text{degree}&4&5&4&22&9&8&26&13&\leq4&9&8&22.
\end{array}
\]
The same exact root identities now give degrees $1,1,-4$ for
$AI-f$, $C_t-\tau$, and $g-V$, respectively.  The decisive small
two-center factor still has degree $5$ with leading form $[I]_5$.

\paragraph{Theorem D.2.}
The direct strong-center polynomial has exact degree $330$, and the
direct two-center polynomial has exact degree $327$.  With
$q_0=\beta J$ and $x_0=q_0+x$, their top forms are
\begin{align*}
 [N_f^{\rm dir}]_{330}
   &=-9\,2^{22}q_0^{160}x_0^{36}Y^{74}i^{18}
       (\eta+\zeta)^{38}(q_0-F)^4,\\
 [N_{ft}^{\rm dir}]_{327}
   &=3\,2^{23}q_0^{156}x_0^{35}Y^{74}i^{19}
       (\eta+\zeta)^{39}(q_0-F)^4.
\end{align*}

\paragraph{Proof.}
Use the same norm factorization and leading-factor classification.
The prefactor contributes $4(26)+8(13)=208$.
The four large factors contribute $4(22)=88$ and the two next factors
contribute $2(9)=18$.  Strong centering leaves two factors of degree $8$,
whereas two-center centering leaves degrees $8$ and $5$.  This gives
$208+88+18+16=330$ and $208+88+18+8+5=327$.

For the top-form calculation replace the earlier abbreviations by
\[
 a_0=q_0^2x_0Y,\quad
 t_0=x_0q_0^4Y^2k_0,\quad m_0=4q_0^9a_0,
\]
retaining $j_0=ik_0^2Y^2$, $r_0=q_0^3(q_0-F)$,
$f_0=a_0j_0$, $v_0=-a_0^2j_0^2r_0$ and
$d_0=a_0^4j_0^2$.  The products
$-576d_0^4m_0^8v_0^4f_0^2t_0^2$ and
$384d_0^4m_0^8v_0^4f_0^2t_0j_0$ give exactly the two displayed
nonzero homogeneous forms. \hfill$\square$

The independent exact ray with all ten witnesses equal to $t$ and
$(\beta,K_0,M_C,M_F)=(2,3,5,7)$ has leading terms
$-9\,2^{220}3^{36}t^{330}$ and
$2^{218}3^{36}t^{327}$.  Both were checked on the supplied literal
104-operation and 106-operation graphs after shared-$C_f$ evaluation.
